% =====================================================================
% Section 6 -- Numerical Experiments and Results (revised)
% Generated by render.py from the experiment outputs; every number below
% is filled in from results/numbers.json. Required packages: booktabs,
% graphicx, amsmath. Figures: figs/*.pdf.
% Cross-references to Sections 3--5 use the numbering of the revised
% manuscript (Eqs. (7)--(23), Propositions 1--4, Corollary 1); replace
% them with \ref{} labels if your source defines them.
% =====================================================================

\section{Numerical Experiments and Results}\label{sec:numerics}

\subsection{Design of the study and evaluation protocol}\label{sec:num-design}

The experiments are organised around five questions, each answered in its own subsection.
\begin{enumerate}
\item[(Q1)] Does the exact dynamic-programming (DP) solution of the base model have the structure asserted by Proposition~2, and do the closed-form expressions of Proposition~3 agree with simulation? (Section~\ref{sec:num-exact})
\item[(Q2)] How large is the value of optimal admission control relative to first-come-first-served (FCFS), and on which model parameters does it depend? (Section~\ref{sec:num-exact})
\item[(Q3)] Does the tabular RVI Q-learning agent of Section~4.2.1 reach the optimum within a finite training budget, and what explains the failures of the configuration used in the original submission? (Section~\ref{sec:num-tabular})
\item[(Q4)] In the extended model of Section~3.5, where arrival regimes are hidden and customers belong to two classes, does a single trained DQN adapt to the regime and to distribution shifts, and how does it compare with classical adaptive and robust controllers? (Section~\ref{sec:num-extended})
\item[(Q5)] Applied to the \emph{trained} DQN, are the explanations of Section~4.3 faithful, stable across seeds, and able to localise where the policy loses reward? (Section~\ref{sec:num-explain})
\end{enumerate}

\paragraph{Exact evaluation wherever possible.} Every policy that depends on the occupancy only is evaluated exactly. Under a deterministic occupancy policy $\pi$ the chain is birth--death with birth rate $\lambda\pi(s)$, so its stationary distribution and $g(\pi)$ follow from the same product form as Proposition~3, whether or not $\pi$ is a threshold rule. This covers FCFS, the best static threshold, the DP policy \emph{and} the greedy policies learned by the tabular agent; Sections~\ref{sec:num-exact}--\ref{sec:num-tabular} therefore contain no simulation noise apart from the explicit check of Proposition~3. In the extended model, any policy that depends on $(s,z,j)$ is evaluated exactly by solving the balance equations of the continuous-time chain on $(s,z)$. Simulation is used only for policies that depend on the regime estimate $\hat\lambda_n$ of Eq.~(7) (DQN, certainty-equivalent DP, plug-in and patched policies) and for the server-failure scenario.

\paragraph{Simulation, uncertainty and reproducibility.} Each simulated estimate uses independent replications of length $T$ after a warm-up of 200 time units; replication $r$ uses the same random-number seed for every controller (common random numbers), so differences between controllers are paired. Unless stated otherwise, intervals are 95\% Student-$t$ intervals across replications; for the DQN they are computed across the five training seeds, each seed being the mean of five replications. Percentage gaps are always relative to the optimum of the same scenario, $(g^*-g)/g^*$, and relative improvements over FCFS are stated as such. The tabular agent is run with 10 seeds and the DQN with 5. All code (Python and NumPy; the DQN is implemented directly in NumPy without a deep-learning framework) and the raw outputs are available at \texttt{[repository URL]}. Table~\ref{tab:hyper} lists all model parameters and hyperparameters; it is the table referred to in Section~4.2.2.

\begin{table}[t]
\centering\small
\caption{Model parameters and hyperparameters used in Section~\ref{sec:numerics}.}\label{tab:hyper}
\begin{tabular}{@{}p{0.27\textwidth}p{0.67\textwidth}@{}}
\toprule
\multicolumn{2}{@{}l}{\emph{Base model (Sections~3.1--3.4)}}\\
System & $c=5$, $K=20$, $\mu=1$, $R=10$, $P=5$, $h=1$; $\lambda\in\{2,\dots,7\}$, offered load $\lambda/(c\mu)\in\{0.4,\dots,1.4\}$; $\Lambda=\lambda+c\mu$\\
RVI Q-learning (Sec.~4.2.1) & reference pair $(s^0,a^0)=(0,1)$; step size $\alpha_k(s,a)=(1+N_k(s,a))^{-\kappa}$, $\kappa=0.6$; $\varepsilon$ linear from $0.6$ to $\varepsilon_{\min}=0.05$ over the first half of training, then constant; $9{,}000$ episodes of $200$ uniformized steps ($1.8\times10^6$ steps, the budget of the original submission); episode start state uniform on $\{0,\dots,K\}$; $Q_0\equiv0$; greedy policy evaluated exactly every $10^4$ steps\\
Legacy R-learning (ablation) & R-learning (Schwartz, 1993) in the configuration of the original submission: constant step size $\alpha=0.1$; gain step $\beta=0.01$, gain updated on greedy steps only; $\varepsilon$ linear from $0.6$ to $0.02$; episodes start in the empty system; same budget\\
\midrule
\multicolumn{2}{@{}l}{\emph{Extended model (Section~3.5)}}\\
System & $c=5$, $K=20$, $\mu=1$, $h=1$; $m=2$ regimes with $(\lambda_1,\lambda_2)=(4,9)$, $\gamma_{12}=\gamma_{21}=0.02$ (mean regime duration 50), $\bar\lambda=6.5$ (offered load 1.3); $q_H=0.3$, $(R_H,P_H)=(20,15)$, $(R_L,P_L)=(10,5)$; window $M=10$ in Eq.~(7)\\
DQN (Sec.~4.2.2) & input $\big(s/K,\ \mathbb 1\{j=H\},\ \min(\hat\lambda,15)/15\big)$; two hidden layers of 64 ReLU units, two outputs; Double-DQN target in Eq.~(11) with infeasible actions masked; Huber loss ($\delta=1$) on rewards scaled by $1/50$; Adam, learning rate $3\times10^{-4}$ decayed linearly to $3\times10^{-5}$; minibatch 128; replay buffer $10^5$; target network synchronised every 1{,}000 updates; 2{,}000 warm-up transitions; $\varepsilon$ linear from 1 to 0.05 over the first 40\% of training; gain smoothing $0.999$ in Eq.~(13); $2\times10^5$ decisions; restart at uniform $(s,z)$ every 2{,}000 decisions; 5 seeds\\
Explanation layer (Sec.~4.3) & 10 rollouts of 3{,}000 time units per seed, 7 for fitting and 3 held out; trees of depth 1--5 and logistic regression; 30 permutations per feature in Eq.~(16); 800 held-out decisions for counterfactuals (Eq.~(17)), $\hat\lambda$ grid step 0.1; 11 bins of width 1 on $\hat\lambda\in[2,13)$ in Eq.~(18)\\
\midrule
Evaluation & exact where available (see text); otherwise replications with common random numbers: Prop.~3 check $10\times20{,}000$; extended-model controllers $10\times10{,}000$ (DQN $5$ seeds $\times\,5$); patch and retraining tests $12\times10{,}000$\\
\bottomrule
\end{tabular}
\end{table}

\subsection{Base model: exact solution, structural checks and the value of control}\label{sec:num-exact}

Table~\ref{tab:base-exact} reports exact performance measures for FCFS and for the optimal policy. The best static threshold (exhaustive search over $T$ with Corollary~1) and relative value iteration on Eq.~(20) return the same policy in every scenario, and their reward rates coincide to within $<<MAX_RVI_EXACT_DIFF>>$; the two are therefore reported as a single ``optimal'' column. The values replace the simulation-based Table~1 of the original submission, in which identical policies carried different reward rates and the best threshold appeared to exceed the DP optimum at $\lambda=5$ purely through simulation noise.

\begin{table}[t]
\centering\small
\caption{Base model: exact performance of FCFS and of the optimal policy ($c=5$, $K=20$, $\mu=1$, $R=10$, $P=5$, $h=1$). The DP solution and the best static threshold coincide. Loss is $(g^*-g_{\rm FCFS})/g^*$.}\label{tab:base-exact}
\begin{tabular}{@{}cc c cccc cccc c@{}}
\toprule
 & & & \multicolumn{4}{c}{FCFS ($T=K$)} & \multicolumn{4}{c}{Optimal (DP $=$ best threshold)} & \\
\cmidrule(lr){4-7}\cmidrule(lr){8-11}
$\lambda$ & $\lambda/c\mu$ & $T^*$ & $P_B$ & $W$ & $U$ & $g$ & $P_B$ & $W$ & $U$ & $g^*$ & Loss (\%)\\
\midrule
<<TAB_BASE_ROWS>>
\bottomrule
\end{tabular}
\end{table}

\paragraph{Direct check of Proposition~2.} Agreement between the DP threshold and the exhaustive-search threshold presupposes that the DP optimum is a threshold rule, so it is only an implementation check. We therefore test the structure directly on the output of relative value iteration, for all $21$ states and all six scenarios: (i) the greedy policy of Eq.~(20) is of threshold form; (ii) the increments $a_s=J^*(s)-J^*(s-1)$ are positive and nondecreasing, i.e.\ $J^*$ is nondecreasing and convex as stated by Lemma~2 and Proposition~2 (the smallest second difference is $<<MIN_D_RANGE>>>0$); and (iii) the threshold equals $s^*$ computed from Eq.~(23). All three properties hold in every scenario.

\paragraph{Check of Proposition~3.} Table~\ref{tab:prop3} compares the closed-form blocking probability, sojourn time and reward rate with replicated discrete-event simulation for FCFS and for $T^*$ in the three scenarios where the two policies differ. The exact value lies inside the 95\% interval for all $<<PROP3_COV>>$ of $<<PROP3_TOT>>$ policy--metric--scenario combinations (including $\lambda\le4$, not shown).

\begin{table}[t]
\centering\small
\caption{Closed form (Proposition~3, Corollary~1) versus simulation (10 replications of 20{,}000 time units; mean $\pm$ 95\% half-width).}\label{tab:prop3}
\begin{tabular}{@{}cl cc cc cc@{}}
\toprule
$\lambda$ & Policy & $P_B$ exact & $P_B$ simulated & $W$ exact & $W$ simulated & $g$ exact & $g$ simulated\\
\midrule
<<TAB_PROP3_ROWS>>
\bottomrule
\end{tabular}
\end{table}

\paragraph{The value of control and what it is attributable to.} Rejecting arrivals before the buffer is full pays off only in heavy traffic. For $\lambda\le4$ the optimal policy is FCFS. For $\lambda=5,6,7$ the optimum rejects from $T^*=14,11,9$ and FCFS loses $<<LOSS5>>$\%, $<<LOSS6>>$\% and $<<LOSS7>>$\% of the optimal reward rate. At $\lambda=7$ the optimal reward rate $<<G7_OPT>>$ exceeds that of FCFS ($<<G7_FCFS>>$) by $<<IMP7>>$\% of the FCFS value. This improvement is a property of the control problem and of the chosen cost parameters; it is obtained by the exact DP solution, and it is \emph{not} a contribution of any learning component. Table~\ref{tab:sens} shows how strongly it depends on the parameters. The threshold falls, and the loss of FCFS grows, when the holding cost is large relative to the value $R+P$ of admitting a customer (the comparison in Eq.~(23)); the loss is small in light traffic unless holding is expensive, and it essentially vanishes when the buffer slack $K-c$ is small: for $(c,K)=(5,10)$ FCFS is within 0.1\% of the optimum at every load. With expensive holding ($h=4$) or a heavy penalty ($P=20$), FCFS can lose more than the entire optimal reward rate. FCFS is thus a weak benchmark in heavy traffic, and the stronger reference throughout the rest of the paper is the exact optimum.

\begin{table}[t]
\centering\small
\caption{Sensitivity of the optimal threshold and of the value of control. Each entry is ``$T^*$; $g^*-g_{\rm FCFS}$'' (exact), varying one parameter at a time from the reference setting; offered load $\rho=\lambda/(c\mu)$.}\label{tab:sens}
\begin{tabular}{@{}lcccc@{}}
\toprule
 & $\rho=0.8$ & $\rho=1.0$ & $\rho=1.2$ & $\rho=1.4$\\
\midrule
<<TAB_SENS_ROWS>>
\bottomrule
\end{tabular}
\end{table}

\subsection{Base model: tabular RVI Q-learning}\label{sec:num-tabular}

\paragraph{Convergence within the original budget.} The agent of Section~4.2.1 was trained with exactly the budget of the original submission ($1.8\times10^6$ uniformized steps). Table~\ref{tab:rl} and Figure~\ref{fig:lc} report the results over 10 seeds. Across all six scenarios the mean gap to the optimum is at most $<<RVI_MEAN_MAX>>$\%, and the worst single seed is $<<RVI_WORST>>$\%. The learned threshold rarely equals $T^*$ exactly at $\lambda=5$ (1 of 10 seeds), yet the gap is small because $g(T)$ is flat near its maximum: every threshold in $<<FLAT5>>$ is within 1\% of $g^*$ at $\lambda=5$, against only $<<FLAT7>>$ at $\lambda=7$. A learned threshold is therefore a reliable diagnostic only where the reward is sensitive to it. The estimate $\hat g=\Lambda Q(0,1)$ of Proposition~4 tracks $g^*$ on average but has a seed-to-seed standard deviation of up to 3.4 at $\lambda=7$; reward rates are therefore reported from exact evaluation of the greedy policy, not from $\hat g$.

\begin{table}[t]
\centering\small
\caption{Tabular learners on the base model, 10 seeds, $1.8\times10^6$ steps, greedy policies evaluated exactly. $T$: occupancy at which the greedy policy first rejects (median [min, max]); ``exact'': seeds with $T=T^*$; gaps $(g^*-g)/g^*$ in \%; $\hat g=\Lambda Q(0,1)$, mean (s.d.).}\label{tab:rl}
\resizebox{\textwidth}{!}{%
\begin{tabular}{@{}cc cccc c ccc@{}}
\toprule
 & & \multicolumn{5}{c}{RVI Q-learning (Section~4.2.1)} & \multicolumn{3}{c}{Legacy R-learning (original submission)}\\
\cmidrule(lr){3-7}\cmidrule(lr){8-10}
$\lambda$ & $T^*$ & $T$ & exact & mean gap & max gap & $\hat g$ & $T$ & mean gap & max gap\\
\midrule
<<TAB_RL_ROWS>>
\bottomrule
\end{tabular}}
\end{table}

\begin{figure}[t]
\centering
\includegraphics[width=\textwidth]{figs/fig_learning_curves.pdf}
\caption{Learning curves for all six traffic scenarios: gap of the greedy policy to $g^*$, evaluated exactly every $10^4$ steps; median (line) and 10--90\% band over 10 seeds. Dashed: FCFS.}\label{fig:lc}
\end{figure}

\paragraph{Why the original configuration failed.} The original submission attributed the learner's failure at $\lambda\ge5$ to under-sampling of the states $9\le s\le19$. Re-running that configuration (constant step sizes, gain updated on greedy steps, episodes starting in the empty system) reproduces the failure: mean gaps of $<<LEG_MEAN6>>$\% and $<<LEG_MEAN7>>$\% at $\lambda=6,7$, and $<<LEG_MAX7>>$\% for the worst seed, which matches the original Table~1. At $\lambda=5$, where the reward is flat around $T^*$, the two configurations perform similarly (the legacy one is slightly better on average); the difference appears where rejection matters most. The diagnostics do not support the under-sampling explanation, which we withdraw. First, the states $9\le s\le19$ are not under-sampled: under the legacy configuration every state--action pair in that range receives at least $<<VIS_LEG_MIN>>$ updates (mean $<<VIS_LEG_MEAN>>$) at $\lambda=7$ (Figure~\ref{fig:visits}). Second, the legacy learner \emph{does} reach the optimal policy: its median seed attains a zero gap after about $<<LEG_STEP_MIN7>>$ thousand steps (Figure~\ref{fig:lc}) and then drifts away as exploration is annealed. The failure is a non-convergence of the update itself, consistent with the violation of the step-size condition~(ii) of Proposition~4 by constant step sizes; it is not a lack of data. The RVI-Q agent has no separately updated gain: the gain is read off $Q(0,1)$, so the reviewer's question of whether the gain update is restricted to greedy transitions applies only to the legacy configuration, where it is.

\begin{figure}[t]
\centering
\includegraphics[width=\textwidth]{figs/fig_visits.pdf}
\caption{Updates per state--action pair at $\lambda=7$, mean over 10 seeds. Shaded: the states $9\le s\le19$ in which rejection is optimal.}\label{fig:visits}
\end{figure}

\paragraph{Sensitivity to step size, exploration and initialisation.} Varying one setting at a time (5 seeds each; Figure~\ref{fig:kappa}) shows that the step-size exponent dominates. At $\lambda=7$ the mean gap is $<<S7_k055_mean>>$\%, $<<S7_base_mean>>$\%, $<<S7_k07_mean>>$\%, $<<S7_k08_mean>>$\% and $<<S7_k10_mean>>$\% for $\kappa=0.55,0.6,0.7,0.8,1$. All five values satisfy Proposition~4, which is asymptotic and does not distinguish between them; within a finite budget, the $1/n$ rate ($\kappa=1$) freezes the policy close to FCFS. The other settings matter much less at $\lambda=7$: $\varepsilon_{\min}\in\{0.01,0.2\}$ gives $<<S7_e001_mean>>$\% and $<<S7_e02_mean>>$\%; episodes starting in the empty system give $<<S7_empty_mean>>$\%; and optimistic initialisation $Q_0\equiv100$ gives $<<S7_q100_mean>>$\%. A five-fold budget brings every seed below $<<S7_x5_max>>$\% at $\lambda=7$ and $<<S5_x5_max>>$\% at $\lambda=5$. Proposition~4 is therefore an asymptotic motivation, not a certificate for the finite training run; the finite-budget evidence is Table~\ref{tab:rl} and Figures~\ref{fig:lc}--\ref{fig:kappa}.

\begin{figure}[t]
\centering
\includegraphics[width=0.5\textwidth]{figs/fig_kappa.pdf}
\caption{Gap to $g^*$ against the step-size exponent $\kappa$ (5 seeds; markers: mean, bars: range).}\label{fig:kappa}
\end{figure}

\paragraph{What explanation adds in the base model.} In the base model the learned policy is a map from 21 occupancy levels to an action. Its natural summary is the occupancy $T$ at which it first rejects, and the counterfactual of Eq.~(17) is the distance to that threshold. As noted in Section~4.3.5, a surrogate adds nothing here, and we do not use one. One observation is nonetheless relevant for auditing. At $\lambda=5$, $6$ and $7$, only $<<MONO_MAIN5>>$, $<<MONO_MAIN6>>$ and $<<MONO_MAIN7>>$ of the 10 learned policies are monotone; the others admit again at some occupancy above the first rejection. Those states are transient under the learned policy, so the reward is unaffected, but a decision rule read off a learned policy must be restricted to its recurrent class. Learners that impose the threshold structure, for example by learning $T$ directly as in Cohen et al.\ (2024) and Weber et al.\ (2024), avoid this issue by construction and are the natural alternative whenever the state is as simple as in the base model.

\subsection{Extended model: hidden regimes, two classes and distribution shift}\label{sec:num-extended}

\paragraph{How much is regime information worth?} Before training any learner, we computed, for several regime settings, the exact optimum of a controller that observes the regime (the regime-aware oracle) and the best static class-specific thresholds $(T_L,T_H)$ (Table~\ref{tab:voi}). In the setting used below, the oracle policy is of threshold form in $s$ for every regime--class pair, which is numerical support for the class-specific threshold structure left open in Section~3.5. The gap between the two bounds what any adaptive controller can gain over a static rule that already exploits the threshold structure. It is small: between $<<VOI_MIN>>$\% and $<<VOI_MAX>>$\% in all settings examined, whereas FCFS loses $<<VOI_FCFS_MIN>>$--$<<VOI_FCFS_MAX>>$\% of the oracle value. Most of the value of control in this model class is therefore captured by structure-exploiting static thresholds, and the scope for adaptivity is limited. We use $(\lambda_1,\lambda_2)=(4,9)$, $\gamma=0.02$ (gap $<<VOI_MAIN>>$\%) for the rest of the section. The oracle there admits high-priority customers until the system is full ($T_H=<<OT_H>>$), and standard customers up to $T_L=<<OT_L_LOW>>$ in the low regime but only $T_L=<<OT_L_HIGH>>$ in the high regime. The regime therefore matters for one class only.

\begin{table}[t]
\centering\small
\caption{Value of regime information (exact). Gaps in \% of the regime-aware optimum. $^\dagger$Setting used in Sections~\ref{sec:num-extended}--\ref{sec:num-explain}.}\label{tab:voi}
\begin{tabular}{@{}cc cc c ccc@{}}
\toprule
$(\lambda_1,\lambda_2)$ & $\gamma$ & Oracle $g$ & Best static $g$ & $(T_L,T_H)$ & Static class-specific & Static class-blind & FCFS\\
\midrule
<<TAB_VOI_ROWS>>
\bottomrule
\end{tabular}
\end{table}

\paragraph{Controllers.} A single DQN (Section~4.2.2) is trained on the training environment $S_0$ and then evaluated without retraining. It is compared with seven controllers, all of which keep the $S_0$ model when the environment shifts.
\begin{itemize}
\item FCFS.
\item The best static class-blind threshold.
\item The best static class-specific thresholds. These are found by exhaustive exact search, and they are the limit that direct threshold-learning methods restricted to static rules can attain.
\item A robust static rule, which maximises the worst-case reward rate over six models with $\lambda_2\in\{7,9,12\}$ and $\mu\in\{0.8,1\}$. This uncertainty set contains two of the shifts below, which favours this baseline.
\item Certainty-equivalent DP. At every arrival it applies the optimal class thresholds of the stationary two-class model with arrival rate $\hat\lambda_n$, i.e.\ DP re-solved on the online estimate. It requires $c$, $\mu$ and the costs, but no regime model.
\item A plug-in regime controller. It classifies the regime by $\hat\lambda_n>6.5$ and applies the oracle's thresholds for that regime; it requires the regime model.
\item The regime-aware oracle, which observes $Z_t$ and serves as an unattainable benchmark.
\end{itemize}
The shifts are: $S_1$, faster switching ($\gamma=0.2$); $S_2$, a demand shock in which the high regime rises from 9 to 12, outside the training range; $S_3$, service degradation to $\mu=0.8$, unknown to every controller; and $S_4$, server failures, in which an on--off process (failure rate 0.02, repair rate 0.1) reduces capacity to $c=3$ while it is off.

\begin{figure}[t]
\centering
\includegraphics[width=0.49\textwidth]{figs/fig_dqn_curve.pdf}
\caption{DQN training on the extended model: simulated reward rate of the greedy policy every $2\times10^4$ decisions (thin: seeds; thick: mean). Reference lines are exact. FCFS ($g=31.24$) lies below the plotted range.}\label{fig:dqn-curve}
\end{figure}

\paragraph{Results.} Training is stable across the five seeds (Figure~\ref{fig:dqn-curve}). The mean reward rate of the greedy policy is $<<DQN_CURVE_100K>>$ after $10^5$ decisions (worst seed $<<DQN_CURVE_100K_MIN>>$) and $<<DQN_CURVE_FINAL>>$ after $2\times10^5$. Table~\ref{tab:scen} reports all controllers in all scenarios. In the training environment the DQN attains $<<DQN0>>$ (seeds $<<DQN0_min>>$--$<<DQN0_max>>$). Paired with the classical controllers on common random numbers, it differs from the best static class-specific rule by $<<D0_cls>>\pm<<D0_cls_hw>>$, from certainty-equivalent DP by $<<D0_ce_dp>>\pm<<D0_ce_dp_hw>>$ and from the plug-in controller by $<<D0_plugin>>\pm<<D0_plugin_hw>>$; the regime-aware optimum is $<<ORC0>>$. Under the shifts the picture is the same. The DQN is statistically indistinguishable from the best classical controller under faster switching ($S_1$), service degradation ($S_3$) and server failures ($S_4$). Under the out-of-range demand shock ($S_2$) it beats the static class-specific rule by $<<D2_cls>>\pm<<D2_cls_hw>>$ but trails certainty-equivalent DP (difference $<<D2_ce_dp>>\pm<<D2_ce_dp_hw>>$). The latter re-solves the stationary model at the observed rate and so extrapolates to $\lambda=12$, which the network never saw. The robust static rule is never better than the nominal static rule, even in the two scenarios its uncertainty set was built to cover.

The evidence therefore does \emph{not} support a claim that the learned policy adapts better than classical methods. The DQN learns, from interaction alone and without access to $\lambda$, $\mu$ or the regime process, a policy that is as good as the best structure-exploiting static rule and close to model-based adaptive control. However, it does not exceed the best classical controller in any scenario, and classical controllers that exploit the threshold structure leave at most $<<VOI_MAIN>>$\% on the table in the training environment. In this model class the case for learning rests on not needing the model, not on better control; whether it rests on more in richer models (heterogeneous service, abandonment, more classes) is left open.

\begin{table}[t]
\centering\small
\caption{Extended model: reward rate of each controller in the training environment $S_0$ and under four shifts (no retraining). Mean $\pm$ 95\% half-width over 10 replications of 10{,}000 time units (DQN: over 5 seeds). Bold: best implementable controller in each column. The oracle observes the hidden regime and is not implementable.}\label{tab:scen}
\resizebox{\textwidth}{!}{%
\begin{tabular}{@{}lccccc@{}}
\toprule
Controller & $S_0$ train & $S_1$ fast switch & $S_2$ demand shock & $S_3$ $\mu=0.8$ & $S_4$ failures\\
\midrule
<<TAB_SCEN_ROWS>>
\bottomrule
\end{tabular}}
\end{table}

\subsection{Explanation layer applied to the trained DQN}\label{sec:num-explain}

This subsection applies the procedure of Section~4.3 to the trained DQN of each seed, not to the DP policy. The surrogates use the agent's own inputs $\phi(o)=(s,j,\hat\lambda)$, which are not deterministic transformations of one another. Rollouts are split by rollout into fitting and test sets (about $<<NFIT>>$ and $<<NTEST>>$ non-forced decisions per seed). Decisions at $s=K$, where rejection is forced, are excluded from all fidelity figures. On average $<<ADMIT_RATE>>$\% of held-out decisions are admissions, and $<<FRAC_NEAR>>$\% are \emph{near-boundary}, meaning a change of occupancy by at most two would reverse the agent's action.

\paragraph{Fidelity (Eq.~(15)).} Table~\ref{tab:fid} reports held-out fidelity overall, on near-boundary decisions, for the decision-relevant stratum (standard customers in the high regime) and on an off-policy grid that covers $s<K$, both classes and $\hat\lambda\in[2,13)$ uniformly. The near-boundary and off-policy figures address the concern that an on-policy average is inflated by states with an obvious action. Depth-4 trees reach $<<FID4>>$ held-out fidelity, but only $<<FID4NB>>$ near the boundary. High-priority decisions are reproduced perfectly by every surrogate, because the agent always admits them. Fitting and held-out fidelity are nearly identical, so the trees do not overfit the rollouts. Seed~0 gives an example of the depth-4 rule: \emph{admit every high-priority customer; admit a standard customer while $s\le9$ if $\hat\lambda\le8.04$ and while $s\le8$ otherwise.}

\begin{table}[t]
\centering\small
\caption{Fidelity of surrogates of the trained DQN: mean [min, max] over 5 seeds. ``Near boundary'': decisions reversed by $|\Delta s|\le2$. ``$z{=}2$, $L$'': standard customers in the high regime. ``Grid'': off-policy grid over $s<K$, $j$, $\hat\lambda\in[2,13)$.}\label{tab:fid}
\begin{tabular}{@{}l c c c c c@{}}
\toprule
Surrogate & Fit & Held out & Near boundary & $z{=}2$, $L$ & Grid\\
\midrule
<<TAB_FID_ROWS>>
\bottomrule
\end{tabular}
\end{table}

\paragraph{Feature importance (Eq.~(16)).} On held-out data, permuting the occupancy reduces the fidelity of the depth-4 tree by $<<PI_S>>$ (range across seeds <<PI_S_R>>), permuting the class by $<<PI_J>>$ (<<PI_J_R>>) and permuting the regime estimate by $<<PI_L>>$ (<<PI_L_R>>). All three inputs matter. How much the policy relies on $\hat\lambda$ varies by a factor of three across seeds, which is itself a finding: equally rewarding policies can weight the regime estimate quite differently. As noted in Section~4.3.2, $s$ and $\hat\lambda$ are correlated through the regime, so these numbers are importances conditional on the observed joint distribution, not causal effects.

\paragraph{Counterfactuals (Eq.~(17)).} For 800 held-out decisions per seed we computed the minimal single-feature change that reverses the agent's action. An occupancy counterfactual exists for $<<CF_S_EXISTS>>$\% of decisions (median $|\Delta s|=<<CF_S_MED>>$; $|\Delta s|\le1$ for $<<CF_S_LE1>>$\%). A $\hat\lambda$ counterfactual exists for $<<CF_L_EXISTS>>$\% of decisions (median $|\Delta\hat\lambda|=<<CF_L_MED>>$). Switching the class reverses $<<CF_J>>$\% of decisions, including every rejection: no rejected standard customer would have been rejected as high priority. A typical statement for a rejected standard customer is: \emph{``rejected at $s=<<EX_S>>$ with $\hat\lambda=<<EX_L>>$; it would have been admitted with one customer fewer in the system ($\Delta s=<<EX_DS>>$), at a regime estimate of $<<EX_DL>>$, or as a high-priority customer.''} Counterfactuals read from the surrogate alone are less reliable. The validity of surrogate-based counterfactuals, i.e.\ the fraction that also reverse the agent's action, is $<<CF_VAL_S>>$\% for occupancy (<<CF_VAL_S_R>>\% across seeds), $<<CF_VAL_L>>$\% for the regime estimate (<<CF_VAL_L_R>>\%) and $<<CF_VAL_J>>$\% for class (<<CF_VAL_J_R>>\%). Their occupancy distance coincides with the agent's in $<<CF_EXACT_S>>$\% of cases. The regime counterfactual is the least reliable, because a shallow tree places its $\hat\lambda$ split at a single value whereas the agent's boundary moves gradually with $\hat\lambda$. A surrogate-based counterfactual should therefore be verified against the agent before it is reported, which is cheap because the agent can be queried.

\paragraph{Threshold profiles and localisation of the loss (Eq.~(18)).} Figure~\ref{fig:profile} plots the threshold profile of the agent, of its surrogate and of two reference controllers for standard customers. High-priority thresholds are at least $<<TA_H_MIN>>$ for every bin and seed, close to the oracle's $T_H=<<OT_H>>$. The cell-wise thresholds of the surrogate agree with those of the agent in $<<PROF_AGREE>>$\% of visit-weighted cells (worst seed $<<PROF_AGREE_MIN>>$\%), with a mean absolute difference of $<<PROF_ABS>>$ customers (worst seed $<<PROF_ABS_MAX>>$). The profile makes the agent's deviation from the reference readable. The plug-in controller switches the standard-class threshold from $<<OT_L_LOW>>$ to $<<OT_L_HIGH>>$ when the regime estimate crosses 6.5. The DQN instead lowers it gradually, from $<<TA_LOW>>$ (seeds <<TA_LOW_R>>) in the lowest bin to $<<TA_HIGH>>$ in the highest, so in the low regime it rejects standard customers well before the reference does. Certainty-equivalent DP also lowers the threshold gradually, but starts higher than the reference (20 in the lowest bin) and ends lower (6 in the highest), because it re-solves the stationary model at each estimate instead of averaging over the two regimes.

\begin{figure}[t]
\centering
\includegraphics[width=0.55\textwidth]{figs/fig_profile.pdf}
\caption{Threshold profile $T_\pi(L,b)$ for standard customers (Eq.~(18)): DQN (thin: seeds; thick: mean), depth-4 tree surrogate (dashed), plug-in regime reference and certainty-equivalent DP. Shaded: bins flagged by the profile comparison in at least three of five seeds.}\label{fig:profile}
\end{figure}

\paragraph{Patch test (Section~4.3.4).} The test checks whether the flagged cells are responsible for the performance gap. The plug-in controller, the best implementable controller in $S_0$, is the reference ($g=<<G_REF>>$ in these replications). A cell $(j,b)$ is flagged when its threshold differs from the reference by at least two and it carries at least 1\% of decisions; in every seed the flagged cells are the standard-class cells at low $\hat\lambda$ (seed~4 also flags three high-$\hat\lambda$ cells). Replacing the agent's decisions by the reference only in those cells recovers $<<PATCH_POOL>>$\% of the pooled gap across seeds. Patches of the same size placed in randomly chosen visited cells recover $<<RAND_POOL>>$\% (Table~\ref{tab:patch}). For the two seeds with the largest gaps the recovered fraction is 87--91\%, and the paired interval excludes zero. For seeds 0 and 2 the gap itself ($\approx0.1$) is within simulation noise, and the test is inconclusive. The profile comparison therefore identifies where the learned policy loses reward without a full re-derivation, but it needs a reference controller. In this experiment the reference is classical and implementable, and it does not require the regime-aware optimum. The claim of the original submission, that the explanation layer detects an under-trained policy \emph{without} any reference, is not supported and has been removed.

\paragraph{Explanation-guided retraining.} We also tested whether the localisation can be turned into better control. Starting from each trained agent, we continued training for $4\times10^4$ decisions in two arms with identical budgets: a guided arm, in which half of the restarts are placed in the flagged cells (the corresponding regime and occupancies between the agent's and the reference's thresholds), and a control arm with ordinary restarts. Relative to the control, the guided arm changes the reward rate by $<<RT_GMC>>$ on average. It is significantly better for $<<RT_POS>>$ seed and significantly worse for $<<RT_NEG>>$, by as much as $<<RT_WORST>>$ (Table~\ref{tab:patch}). Inspecting the retrained profiles shows why. Guided restarts raise the low-regime thresholds as intended, but they over-represent those states in the replay buffer and distort the policy in cells that were not flagged. The explanation localises the loss but does not by itself repair it, and naive targeted retraining can make the policy worse.

\begin{table}[t]
\centering\small
\caption{Patch test and explanation-guided retraining per seed (12 paired replications of 10{,}000 time units). Gap: $g_{\rm ref}-g_{\rm agent}$. Flagged cells and their share of decisions. Gains are paired differences to the unpatched agent (mean $\pm$ 95\% half-width); ``recovered'' is patch gain over gap.}\label{tab:patch}
\resizebox{\textwidth}{!}{%
\begin{tabular}{@{}c c c c c c c@{}}
\toprule
Seed & Gap & Flagged cells & Patch gain & Recovered (\%) & Random patch gain & Guided $-$ control retraining\\
\midrule
<<TAB_PATCH_ROWS>>
\bottomrule
\end{tabular}}
\end{table}

\paragraph{Stability across seeds (Eq.~(19)).} On a common sample of 5{,}000 held-out states, the depth-4 surrogates of different seeds agree on $<<STAB_S>>$\% of decisions on average (lowest pair $<<STAB_S_MIN>>$\%); the agents themselves agree on $<<STAB_A>>$\% (lowest pair $<<STAB_A_MIN>>$\%). The visit-weighted across-seed standard deviation of the implied thresholds is $<<THR_SD_A>>$ customers for the agents and $<<THR_SD_S>>$ for the surrogates. It is largest in the lowest $\hat\lambda$ bin ($<<THR_SD_LOWBIN>>$ for the agents), where decisions are rarest. Explanations are therefore as stable as the policies they explain, but not more: about one decision in nine differs between independently trained agents with statistically indistinguishable reward rates. An explanation of one trained agent should not be read as an explanation of the learning method.

\subsection{Summary of the evidence}\label{sec:num-summary}

\begin{enumerate}
\item \emph{Structure and closed forms (Q1).} On the base model the exact DP solution has the threshold structure of Proposition~2 and a nondecreasing convex relative cost in every scenario, verified state by state. Proposition~3 agrees with simulation in all $<<PROP3_TOT>>$ comparisons.
\item \emph{Value of control (Q2).} Optimal admission control improves on FCFS by up to $<<IMP7>>$\% at $\lambda=7$ in the reference setting. The improvement comes from the control problem, not from learning, and it depends strongly on $h/(R+P)$, the load and the buffer slack $K-c$ (Table~\ref{tab:sens}).
\item \emph{Tabular learning (Q3).} With the step sizes required by Proposition~4 and the original budget, tabular RVI Q-learning is within $<<RVI_MEAN_MAX>>$\% of the optimum on average in every scenario. The failure reported in the original submission is reproduced by its configuration and is caused by non-convergent step sizes, not by under-sampling.
\item \emph{Adaptivity (Q4).} In the extended model a single DQN, trained without access to the model, matches the best static class-specific thresholds and is close to model-based adaptive control, in the training environment and under four shifts. It does not outperform them, and regime information is worth at most $<<VOI_MAX>>$\% over a static class-specific rule in the settings examined.
\item \emph{Explanation (Q5).} On the trained DQN, depth-4 surrogates are faithful on held-out data ($<<FID4>>$), less so near the decision boundary ($<<FID4NB>>$). Class and occupancy counterfactuals are reliable when verified against the agent. Profile comparison with a classical reference localises most of the performance gap ($<<PATCH_POOL>>$\% recovered against $<<RAND_POOL>>$\% for random patches). Naive explanation-guided retraining did not improve the policy.
\end{enumerate}
These results do not establish that the explanations are useful to human operators, which would require a user study. Nor do they cover policy-gradient agents (PPO, SAC), non-exponential service, abandonment, or more than two classes.
